\chapter{Trisomy 18}
\index{Trisomy 18}

\section*{Definition and Overview}
Trisomy 18, historically and commonly referred to as Edwards syndrome, is a severe genetic disorder characterized by the presence of a complete or partial extra copy of chromosome 18. First described in the medical literature in 1960 by British geneticist John Hilton Edwards and his colleagues, it stands as the second most common autosomal trisomy among live-born infants, following Trisomy 21 (Down syndrome). Human somatic cells typically contain 23 pairs of chromosomes, comprising 22 pairs of autosomes and one pair of sex chromosomes, for a total of 46. In individuals with Trisomy 18, this genetic blueprint is altered by a third copy of chromosome 18, bringing the total chromosome count to 47. This excess genetic material profoundly disrupts normal embryonic and fetal development, leading to a high rate of spontaneous miscarriage and stillbirth. For those infants who survive to birth, the condition manifests as a complex clinical syndrome featuring severe intrauterine growth restriction, multi-system congenital malformations, structural heart defects, developmental delays, and a significantly shortened lifespan.

Understanding Trisomy 18 is of paramount importance to clinical geneticists, obstetricians, neonatologists, and pediatric specialists. It presents a complex array of clinical, developmental, and ethical challenges that span from initial prenatal screening and genetic counseling to postnatal medical management. Historically, the diagnosis was considered uniformly lethal, with medical care focused strictly on non-intervention. However, contemporary management increasingly incorporates a shared decision-making model between clinicians and families. This model balances palliative comfort care with targeted surgical and medical interventions to improve survival and quality of life for select children. The clinical, emotional, and ethical dimensions of Trisomy 18 underscore its significance in modern medicine, illustrating the evolution of neonatal care, parental advocacy, and the integration of pediatric palliative care networks.

\section*{Epidemiology}
The epidemiological landscape of Trisomy 18 is defined by its relation to maternal age, high rates of fetal loss, and a skewed sex ratio at birth. It is estimated to occur in approximately 1 in 5,000 to 1 in 8,000 live births worldwide. This frequency, however, underrepresents the true incidence of the disorder at conception. The vast majority of pregnancies affected by Trisomy 18 do not survive to term. It is estimated that between 60\% and 70\% of fetuses diagnosed with the condition prenatally result in spontaneous abortion, miscarriage, or stillbirth during the second or third trimesters. Consequently, the incidence of Trisomy 18 in early pregnancy—typically evaluated during first-trimester screening at 10 to 12 weeks of gestation—is substantially higher, estimated to be around 1 in 2,500 conceptions.

The risk of carrying a fetus with Trisomy 18 is strongly associated with advancing maternal age. The probability of meiotic error increases as maternal age rises, with a notable acceleration in risk for women aged 35 and older. This relationship is due to the increased frequency of maternal meiotic nondisjunction in older oocytes. Despite this risk curve, it is important to recognize that because total birth rates are higher among younger women, a substantial number of babies with Trisomy 18 are born to mothers under the age of 35.

A striking epidemiological feature of Trisomy 18 is the sex distribution observed at birth. There is a pronounced female-to-male ratio of approximately 3:1 among live-born infants. Interestingly, studies utilizing prenatal diagnosis have shown that the sex ratio at the time of early amniocentesis or chorionic villus sampling is closer to 1:1. This finding indicates that male fetuses with Trisomy 18 experience a significantly higher rate of in-utero mortality compared to female fetuses. The precise genetic or physiological factors that render male fetuses more vulnerable to intrauterine demise remain a subject of active medical research.

\section*{Aetiology and Pathophysiology}
The underlying cause of Trisomy 18 is the presence of three copies of chromosome 18 in somatic cells, which disrupts the dosage-sensitive pathways of embryonic gene expression. The etiology and pathophysiological mechanisms of the condition can be broken down as follows:
\begin{itemize}
    \item \textbf{Meiotic Nondisjunction (Full Trisomy 18)}: This is the most common cause, accounting for approximately 90\% to 95\% of all diagnosed cases. Nondisjunction refers to the failure of homologous chromosomes or sister chromatids to separate during cell division. In full Trisomy 18, this error occurs during meiosis, the specialized cell division that produces eggs and sperm. When an abnormal gamete containing an extra chromosome 18 fuses with a normal gamete, it forms a zygote with 47 chromosomes. Molecular studies have demonstrated that in approximately 90\% of these cases, the nondisjunction event occurs during maternal oogenesis, with errors in maternal meiosis II being the most frequent source.
    \item \textbf{Mosaic Trisomy 18}: This form occurs in approximately 5\% of patients. It is characterized by the presence of two distinct cell lines in the individual's tissues: one cell line with the normal complement of 46 chromosomes and another containing the extra chromosome 18. Mosaicism typically arises from a post-zygotic mitotic nondisjunction event in an initially normal embryo, or from ``trisomy rescue,'' where a cell in a trisomic embryo loses the extra chromosome during division. The clinical severity of mosaic Trisomy 18 is highly variable and unpredictable, as it is determined by the proportion of trisomic cells and their specific distribution throughout the body's organ systems. Some mosaic individuals exhibit milder dysmorphic features and better cognitive development, while others are as severely affected as those with full trisomy.
    \item \textbf{Partial Trisomy 18}: Representing less than 1\% of cases, partial trisomy occurs when only a portion of chromosome 18 is duplicated. This is usually the result of a structural chromosome rearrangement, such as a balanced translocation in one parent that is inherited in an unbalanced form, or a de novo duplication. The clinical presentation of partial Trisomy 18 depends on which segment of the chromosome is duplicated and the size of the duplicated region. Research suggests that duplication of the distal long arm (specifically the region from 18q21 to 18qter) is critical for producing the classic clinical phenotype of Edwards syndrome.
    \item \textbf{Gene Dosage Effect and Pathophysiology}: Chromosome 18 contains approximately 300 to 400 genes. The presence of three copies instead of two leads to a 1.5-fold increase in the expression of these genes, disrupting delicate gene networks that coordinate organogenesis. The overexpression of genes like \textit{NEDD4} (which regulates protein degradation and cell growth), \textit{ASB15} (involved in skeletal muscle differentiation), and various transcription factors leads to cellular dysfunction, impaired mitotic cycles, and dysregulated apoptosis. This disruption alters embryonic development at critical stages, leading to the characteristic structural malformations.
    \item \textbf{Risk Factors}: Advanced maternal age is the only well-established risk factor for meiotic nondisjunction in Trisomy 18. The age-related degradation of cohesin complexes, which hold sister chromatids together during the long meiotic arrest of oocytes, increases the risk of segregation errors. There are no known environmental, occupational, or lifestyle factors that cause or prevent Trisomy 18.
\end{itemize}

\section*{Clinical Features}
The clinical presentation of Trisomy 18 is characterized by a high burden of structural anomalies, distinct craniofacial features, and multi-system dysfunction. These findings are often recognized during prenatal ultrasounds or immediately upon birth.
\begin{itemize}
    \item \textbf{Prenatal Presentation}: During gestation, fetuses with Trisomy 18 display severe, early-onset intrauterine growth restriction (IUGR). Ultrasound findings include a single umbilical artery, choroid plexus cysts in the brain, a strawberry-shaped skull, polyhydramnios (excess amniotic fluid due to poor swallowing), or oligohydramnios (reduced fluid due to renal anomalies). Decreased fetal movement and an abnormal biophysical profile are also common.
    \item \textbf{Craniofacial Anomalies}: The head and face exhibit a highly characteristic appearance. Newborns typically present with microcephaly (an abnormally small head) and a prominent occiput (the back of the skull). The palpebral fissures (eye openings) are narrow, and microphthalmia (small eyes) may be present. The ears are low-set, small, and folded (often referred to as ``faun-like'' ears). The jaw is severely underdeveloped (micrognathia), resulting in a small mouth. Cleft lip and cleft palate are seen in a minority of cases.
    \item \textbf{Skeletal and Limb Features}: The limbs show classic anomalies that strongly suggest the diagnosis. The hands are typically clenched, with the index finger overlapping the middle finger, and the fifth finger overlapping the fourth finger. The fingernails and toenails are hypoplastic (underdeveloped) or dysplastic. The feet have a characteristic ``rocker-bottom'' appearance, featuring a prominent heel bone (calcaneus) and a rounded sole. A short sternum (breastbone), narrow pelvis, and joint contractures are also common.
    \item \textbf{Cardiovascular Malformations}: Congenital heart disease is present in more than 90\% of infants with Trisomy 18. The most common defects are ventricular septal defects (VSD) and atrial septal defects (ASD), often accompanied by a patent ductus arteriosus (PDA). More complex cardiac anomalies include coarctation of the aorta, tetralogy of Fallot, transposition of the great arteries, hypoplastic left heart syndrome, and polyvalvular disease.
    \item \textbf{Gastrointestinal and Abdominal Anomalies}: Malformations of the gastrointestinal tract include esophageal atresia, frequently with a tracheoesophageal fistula (TEF). This combination prevents swallowing and carries a high risk of aspiration. Omphalocele, where abdominal contents protrude through the umbilical ring covered by a thin membrane, is another classic feature. Diaphragmatic hernia, umbilical hernias, and inguinal hernias are also frequently observed.
    \item \textbf{Genitourinary Malformations}: Cryptorchidism (undescended testes) is present in almost all male infants, along with hypoplasia of the scrotum. Female infants often have hypoplastic labia majora. The most common renal anomaly is a horseshoe kidney, where the lower poles of both kidneys are fused across the midline. Other anomalies include renal dysplasia, double ureters, and hydronephrosis.
    \item \textbf{Neurological Manifestations}: Severe developmental delay and intellectual disability are universal in survivors. Neonates often display marked hypotonia (reduced muscle tone) in the immediate postnatal period, which typically transitions to hypertonia and spasticity over time. The sucking and swallowing reflexes are poorly coordinated or absent, necessitating enteral feeding. Central apnea and seizures are also common.
\end{itemize}

\section*{Differential Diagnosis}
Given the wide range of congenital anomalies associated with Trisomy 18, a careful differential diagnosis is necessary to distinguish it from other genetic syndromes and malformation associations:
\begin{itemize}
    \item \textbf{Trisomy 13 (Patau Syndrome)}: This autosomal trisomy presents with overlapping features like microcephaly, congenital heart disease, and severe developmental delays. However, Trisomy 13 is distinguished by midline anomalies, including holoprosencephaly, cleft lip and palate, microphthalmia, and postaxial polydactyly (extra digits). While rocker-bottom feet can occur in Trisomy 13, they are less characteristic than in Trisomy 18.
    \item \textbf{Trisomy 21 (Down Syndrome)}: While also presenting with hypotonia and cardiac defects, Trisomy 21 is characterized by a distinct phenotype including upslanting palpebral fissures, a flat facial profile, a single transverse palmar crease, Brushfield spots in the iris, and a lack of the severe skeletal dysmorphism (e.g., overlapping fingers, rocker-bottom feet) seen in Trisomy 18.
    \item \textbf{Pena-Shokeir Syndrome (Type I)}: An autosomal recessive disorder that closely mimics Trisomy 18, presenting with intrauterine growth restriction, multiple joint contractures (arthrogryposis), micrognathia, low-set ears, and pulmonary hypoplasia. It is distinguished from Trisomy 18 by a normal karyotype.
    \item \textbf{Smith-Lemli-Opitz Syndrome}: An autosomal recessive metabolic disorder of cholesterol biosynthesis caused by mutations in the \textit{DHCR7} gene. Symptoms include microcephaly, developmental delay, syndactyly of the second and third toes, cleft palate, and genital anomalies. Diagnosis is confirmed by demonstrating elevated levels of 7-dehydrocholesterol in the blood.
    \item \textbf{CHARGE Syndrome}: A genetic disorder characterized by Coloboma, Heart defects, Atresia choanae, Restriction of growth, Genital anomalies, and Ear abnormalities. It is caused by mutations in the \textit{CHD7} gene. While it shares several structural anomalies with Trisomy 18, it is cytogenetically normal and does not present with the typical limb anomalies.
    \item \textbf{VACTERL Association}: A non-random clustering of congenital anomalies including Vertebral defects, Anal atresia, Cardiac defects, Tracheo-esophageal fistula, Renal anomalies, and Limb abnormalities. Unlike Trisomy 18, VACTERL association is typically not associated with chromosomal anomalies, and affected infants do not have the characteristic facial dysmorphism or profound intellectual impairment.
\end{itemize}

\section*{When to Seek Medical Care}
For parents expecting an infant with a prenatal diagnosis of Trisomy 18, close communication with a maternal-fetal medicine specialist, neonatologist, and pediatric palliative care physician is essential. A detailed birth plan, outlining preferred levels of intervention, should be established prior to delivery.

During pregnancy, immediate medical attention should be sought if there is a sudden decrease in fetal movement, vaginal bleeding, amniotic fluid leakage, or signs of preterm contractions.

Postnatally, infants with Trisomy 18 are highly vulnerable due to structural anomalies and immature physiological systems. Parents and caregivers must seek immediate emergency medical care if the infant displays any of the following warning signs:
\begin{itemize}
    \item Cyanosis (a blue, purple, or gray discoloration of the skin, lips, or tongue), indicating insufficient blood oxygen levels.
    \item Signs of severe respiratory distress, including rapid breathing, flaring of the nostrils, grunting, or deep indentations of the chest wall (intercostal and subcostal retractions).
    \item Apnea (pauses in breathing lasting more than 15 to 20 seconds, or shorter pauses accompanied by limpness or color change).
    \item Feeding difficulties accompanied by coughing, choking, or vomiting, which can indicate aspiration.
    \item A sudden change in consciousness, extreme lethargy, or repetitive, rhythmic movements of the limbs suggesting seizures.
    \item Fever (rectal temperature above 38 degrees Celsius or 100.4 degrees Fahrenheit) or hypothermia (rectal temperature below 36 degrees Celsius or 96.8 degrees Fahrenheit), especially when accompanied by poor feeding.
\end{itemize}

In the event of an acute, life-threatening emergency, caregivers should call their local emergency number immediately, such as \textbf{local emergency services} in many countries, \textbf{911} in the United States, or \textbf{999} in the United Kingdom. If the family has chosen a palliative path, they should follow the designated contact protocols provided by their palliative care provider.

\section*{Diagnosis and Investigations}
The diagnosis of Trisomy 18 involves prenatal screening, diagnostic genetic testing, and postnatal assessments to evaluate the severity of organ malformations.
\begin{itemize}
    \item \textbf{Prenatal Screening}: Non-invasive prenatal testing (NIPT) utilizing cell-free fetal DNA (cfDNA) in maternal blood can detect Trisomy 18 with high accuracy (sensitivity greater than 99\%) starting at 10 weeks of gestation. First-trimester combined screening (ultrasound measurement of nuchal translucency plus maternal serum PAPP-A and $\beta$-hCG) and second-trimester quadruple screening are also effective screening tools.
    \item \textbf{Prenatal Diagnostic Testing}: Screening results must be confirmed with diagnostic testing. Chorionic villus sampling (CVS) is performed between 10 and 14 weeks of gestation, and amniocentesis is performed after 15 weeks. These procedures allow for cytogenetic analysis of fetal cells, providing a definitive diagnosis.
    \item \textbf{Postnatal Diagnosis}: If the diagnosis is not made prenatally, it is suspected based on the infant's physical features and confirmed postnatally. A rapid fluorescent in situ hybridization (FISH) test provides preliminary results, followed by a conventional karyotype analysis of peripheral blood lymphocytes to confirm the diagnosis and identify the cytogenetic subtype (full, mosaic, or partial).
    \item \textbf{Post-Diagnostic Investigations}: Once a diagnosis is confirmed, several investigations are performed to guide management:
    \begin{itemize}
        \item \textit{Echocardiography}: A detailed cardiac ultrasound is essential to identify structural heart defects and assess ventricular function.
        \item \textit{Renal Ultrasound}: Performed to screen for horseshoe kidney, hydronephrosis, or other structural genitourinary anomalies.
        \item \textit{Brain Imaging}: Cranial ultrasound or magnetic resonance imaging (MRI) is used to evaluate central nervous system malformations.
        \item \textit{Gastrointestinal Studies}: Contrast studies or ultrasound may be indicated if esophageal atresia or malrotation is suspected.
        \item \textit{Sensory Testing}: Hearing evaluations (brainstem auditory evoked response) and ophthalmology examinations are conducted to identify sensory deficits.
    \end{itemize}
\end{itemize}

\section*{Management}
The management of Trisomy 18 is highly individualized, requiring a multidisciplinary approach that includes geneticists, neonatologists, pediatric cardiologists, palliative care specialists, therapists, and social workers. Management strategies have shifted toward a family-centered approach, balancing comfort care with targeted medical and surgical interventions.
\begin{itemize}
    \item \textbf{Perinatal Palliative Care}: For many families, the primary focus of care is comfort and maximizing the quality of life. Palliative care teams work with families to manage symptoms, provide warmth, control pain, and support feeding. Comfort care can be provided in the hospital, a pediatric hospice facility, or at home.
    \item \textbf{Feeding and Nutritional Support}: Due to poor coordination of sucking and swallowing, enteral feeding is usually required. Short-term support is provided via a nasogastric (NG) or orogastric (OG) tube. In infants who survive the neonatal period and show stable growth, surgical placement of a gastrostomy tube (G-tube) may be performed to facilitate long-term feeding and reduce the risk of aspiration.
    \item \textbf{Respiratory Interventions}: Management of respiratory symptoms depends on the goals of care. Non-invasive support includes supplemental oxygen, high-flow nasal cannula, or continuous positive airway pressure (CPAP) to manage central or obstructive apnea. In some cases, mechanical ventilation or tracheostomy may be considered for long-term airway management.
    \item \textbf{Cardiovascular Interventions}: Standard medical management of congestive heart failure includes diuretics (e.g., furosemide) to reduce fluid retention and afterload reducers. While cardiac surgery was historically not offered to infants with Trisomy 18, recent clinical evidence indicates that select infants who undergo surgical repair of congenital heart defects (e.g., VSD repair, ductal ligation) experience improved survival and reduction in heart failure symptoms.
    \item \textbf{Neurological and Supportive Care}: Seizures are treated with anticonvulsant medications. Physical, occupational, and speech therapy are introduced early to help manage muscle tone, prevent joint contractures, and support developmental milestones.
    \item \textbf{Psychosocial Support}: Emotional and psychological support for the family is a cornerstone of management. Social workers, counselors, and spiritual care providers assist families in navigating the emotional challenges, and genetic counseling is provided to discuss future family planning.
\end{itemize}

\section*{Complications}
Infants with Trisomy 18 face a high incidence of complications due to their underlying chromosomal abnormality. These complications can affect multiple organ systems and are major contributors to morbidity and mortality.
\begin{itemize}
    \item \textbf{Cardiovascular Failure}: Structural heart defects, particularly large VSDs and PDAs, cause severe left-to-right shunts, leading to progressive congestive heart failure. Chronic volume overload can result in pulmonary hypertension, which carries a poor prognosis and limits therapeutic options.
    \item \textbf{Respiratory Failure}: Central apnea, caused by dysfunction of the brainstem respiratory center, and obstructive apnea, due to micrognathia and retroposition of the tongue, are common. Aspiration of secretions or stomach contents can cause recurrent aspiration pneumonia, a frequent cause of death.
    \item \textbf{Gastrointestinal and Nutritional Complications}: Severe feeding difficulties can lead to gastroesophageal reflux disease (GERD) and chronic failure to thrive. Intestinal obstruction may occur due to malrotation or volvulus. Esophageal atresia and tracheoesophageal fistula require complex surgical repair, which carries significant risks in these fragile infants.
    \item \textbf{Renal Complications}: Structural kidney anomalies, such as a horseshoe kidney, increase susceptibility to urinary tract infections and pyelonephritis. These infections can lead to renal scarring and chronic kidney disease.
    \item \textbf{Oncological Risks}: Children with Trisomy 18 who survive past the first year of life have a significantly increased risk of developing embryonal tumors, specifically Wilms' tumor (nephroblastoma) and hepatoblastoma. Regular screening with abdominal ultrasound and serum alpha-fetoprotein (AFP) monitoring is often recommended.
    \item \textbf{Musculoskeletal and Neurological Issues}: Severe scoliosis can develop as the child grows, further compromising respiratory function. Refractory seizures, spasticity, and joint contractures require ongoing management and can significantly affect the child's comfort and mobility.
\end{itemize}

\section*{Prognosis and Outcomes}
The prognosis for infants with Trisomy 18 is generally characterized by high mortality rates prenatally and in the early neonatal period. The median survival for live-born infants is approximately 3 to 14 days, with only about 50\% of infants surviving past the first week of life. The primary causes of neonatal death include central apnea, progressive cardiac failure, and aspiration pneumonia.

Despite the high mortality rate, survival outcomes are variable. Approximately 20\% to 30\% of live-born infants survive to one month, and approximately 5\% to 10\% survive past the first year of life. Factors that favor survival include female sex, higher birth weight, and the absence of major congenital heart defects or diaphragmatic hernias. Furthermore, individuals with mosaic or partial Trisomy 18 typically have a more favorable prognosis and a longer life expectancy compared to those with full Trisomy 18, though their outcomes are highly dependent on the proportion and tissue distribution of the trisomic cells.

For the small cohort of children who survive past infancy, development is marked by profound intellectual disability and motor delays. Most survivors do not achieve independent ambulation or verbal communication. However, many are able to interact with their families, show recognition of parents and caregivers, smile, make vocalizations, and achieve basic developmental milestones like sitting with support or rolling over. These children require comprehensive, lifelong medical care and developmental support.

\section*{Prevention and Self-Care}
As Trisomy 18 is a genetic disorder caused by a random error in chromosome segregation (nondisjunction) during gamete formation, there are no specific prevention strategies or lifestyle interventions that can prevent the condition from occurring. It is not caused by any paternal or maternal behavior.

However, several steps can be taken for risk assessment, family preparation, and home care:
\begin{itemize}
    \item \textbf{Genetic Counseling}: Meeting with a genetic counselor is highly recommended for parents who have had a pregnancy affected by Trisomy 18. For families where the child has full Trisomy 18, the recurrence risk in a subsequent pregnancy is low, typically estimated at less than 1\%, or slightly higher depending on maternal age. However, if the child has a partial translocation Trisomy 18, parental karyotyping is performed to determine if one parent carries a balanced translocation. If a parent is a carrier, the risk of recurrence is significantly higher, and the counselor can outline specific inheritance patterns.
    \item \textbf{Prenatal Options for Future Pregnancies}: Parents who wish to avoid a recurrence may consider preimplantation genetic testing (PGT) combined with in vitro fertilization (IVF). This allows for the selection of embryos that are cytogenetically normal before implantation. Alternatively, early prenatal diagnostic testing (CVS or amniocentesis) can be utilized in future pregnancies to provide early diagnostic clarity.
    \item \textbf{Comprehensive Care Planning}: For families choosing to continue a pregnancy affected by Trisomy 18, developing a clear perinatal palliative care plan is crucial. This involves coordinating with obstetricians, neonatologists, and palliative care specialists to document the family's choices regarding resuscitation, intensive care interventions, and comfort measures prior to delivery.
    \item \textbf{Home Care and Self-Care for Caregivers}: Caring for a child with Trisomy 18 at home requires significant physical and emotional energy. Self-care strategies and practical steps for caregivers include:
    \begin{itemize}
        \item Utilizing home nursing support, respite care services, and physical therapy resources to manage daily care requirements.
        \item Learning techniques for safe enteral feeding (G-tube care, positioning) and airway management (suctioning, monitoring).
        \item Participating in support groups and counseling to process grief, cope with chronic stress, and connect with other families facing similar diagnoses.
    \end{itemize}
\end{itemize}

\section*{Sources and Further Reading}
\begin{itemize}
    \item National Organization for Rare Disorders (NORD). Trisomy 18. URL: \url{https://rarediseases.org/rare-diseases/trisomy-18-syndrome/}
    \item Support Organization for Trisomy 18, 13 and Related Disorders (SOFT). Information and Family Support Services. URL: \url{https://trisomy.org}
    \item Trisomy 18 Foundation. Resources, Education, and Research Funding. URL: \url{https://www.trisomy18.org}
    \item MedlinePlus (US National Library of Medicine). Genetics: Trisomy 18. URL: \url{https://medlineplus.gov/genetics/condition/trisomy-18/}
    \item Centers for Disease Control and Prevention (CDC). Facts about Trisomy 18. URL: \url{https://www.cdc.gov/ncbddd/birthdefects/trisomy18.html}
    \item Orphanet. The Portal for Rare Diseases and Orphan Drugs: Edwards Syndrome. URL: \url{https://www.orpha.net}
\end{itemize}